\documentclass[11pt,a4paper]{article}
			\usepackage[french]{babel}
					
				\usepackage{pifont}  
				\usepackage[utf8x]{inputenc}
				\usepackage[T1]{fontenc} 
				\usepackage{lmodern}			
				\usepackage{fancyhdr}
				\usepackage{textcomp}
				\usepackage{makeidx}
				\usepackage{tabularx}
				\usepackage{multicol}
				\usepackage{multirow}
				\usepackage{longtable}
				\usepackage{color}
				\usepackage{soul}
				\usepackage{boxedminipage}
				\usepackage{shadow}
				\usepackage{framed}			
				\usepackage{array}
				\usepackage{url}
				\usepackage{ragged2e}
				\usepackage{fancybox}
				\newcommand{\cadretitre}[2]{
				  \vspace*{0.8\baselineskip}
				  \begin{center}%
				  \boxput*(0,1){%
					%\colorbox{white}{\Large\textbf{\ #1\ }}%
				  }%
				  {%
					\setlength{\fboxsep}{10pt}%
				    \Ovalbox{\begin{minipage}{.8\linewidth}\begin{center}\Large\sffamily{#2}\end{center}\end{minipage}}}%
				  \end{center}
				  \vspace*{2\baselineskip}
				  }
			
			\makeatletter
			\def\@seccntformat#1{\protect\makebox[0pt][r]{\csname the#1\endcsname\quad}}
			\makeatother

				% Permet d'afficher qqchose à une positin absolue
				\usepackage[absolute]{textpos}
				\setlength{\TPHorizModule}{1cm}
				\setlength{\TPVertModule}{\TPHorizModule}
	
				\usepackage[titles]{tocloft}
				\setlength{\cftbeforesecskip}{0.5ex}
				\setlength{\cftbeforesubsecskip}{0.2ex}
				\addto\captionsfrench{\renewcommand\contentsname{}}
				
				\usepackage[font=scriptsize]{caption}
				
				\usepackage{listings}
\lstdefinestyle{lstverb}
  {
    basicstyle=\footnotesize,
    frameround=tttt, frame=trbl, framerule=0pt, rulecolor=\color{gray},
    lineskip=-1pt,   % pour rapprocher les lignes
    flexiblecolumns, escapechar=\\,
    tabsize=4, extendedchars=true
  }
\lstnewenvironment{Java}[1][]{\lstset{style=lstverb,language=java,#1}}{}
				\ifx\pdfoutput\undefined
					\usepackage{graphicx}
				\else
					\usepackage[pdftex]{graphicx}
				\fi
				\usepackage[a4paper, hyperfigures=true, colorlinks, linkcolor=black, citecolor=blue,urlcolor=blue, pagebackref=true, bookmarks=true, bookmarksopen=true,bookmarksnumbered=true,
                pdfauthor={}, pdftitle={TD Tri}, pdfkeywords={TD Tri, },pdfpagemode=UseOutlines,pdfpagetransition=Dissolve,nesting=true,
				backref, pdffitwindow=true, bookmarksnumbered=true]{hyperref}
				\usepackage{supertabular}
				\usepackage[table]{xcolor}
				\usepackage{url}
				\usepackage{caption} 
				\setlength{\parskip}{1.3ex plus 0.2ex minus 0.2ex}
				\setlength{\parindent}{0pt}
				
				\makeatletter
				\def\url@leostyle{ \@ifundefined{selectfont}{\def\UrlFont{\sf}}{\def\UrlFont{\footnotesize\ttfamily}}}
				\makeatother
				\urlstyle{leo}
				
				\definecolor{examplecolor}{rgb}{0.156,0.333,0.443}
				\definecolor{definitioncolor}{rgb}{0.709,0.784,0.454}
				\definecolor{exercisecolor}{rgb}{0.49,0.639,0}
				\definecolor{hintcolor}{rgb}{0.941,0.674,0.196}
				\definecolor{tableHeadercolor}{rgb}{0.709,0.784,0.454}
				\definecolor{tablerowAltcolor}{rgb}{.866,.905,.737}
				\definecolor{tablerowAlt2color}{rgb}{.968,.976,.933}
				\definecolor{verylightgray}{rgb}{0.98,0.98,0.98}
				
				\newenvironment{fshaded}{
				\def\FrameCommand{\fcolorbox{framecolor}{shadecolor}}
				\MakeFramed {\FrameRestore}}
				{\endMakeFramed}
				
				\newenvironment{fexample}[1][]{\definecolor{shadecolor}{rgb}{.913,.913,.913}
				\definecolor{framecolor}{rgb}{.156,.333,.443}
				\begin{fshaded}}{\end{fshaded}} 
				
				\newenvironment{fdefinition}{\definecolor{shadecolor}{rgb}{.913,.913,.913}
				\definecolor{framecolor}{rgb}{.709,.784,.454}
				\begin{fshaded}}{\end{fshaded}}
				
				\newenvironment{fexercise}{\definecolor{shadecolor}{rgb}{.913,.913,.913}
				\definecolor{framecolor}{rgb}{.49,.639,0}
				\begin{fshaded}}{\end{fshaded}}
				
				\newenvironment{fhint}{\definecolor{shadecolor}{rgb}{.913,.913,.913}
				\definecolor{framecolor}{rgb}{.941,.674,.196}
				\begin{fshaded}}{\end{fshaded}}	
				
				\newcommand{\PreserveBackslash}[1]{
				\let\temp=\\#1\let\\=\temp
				}
				\let\PBS=\PreserveBackslash
				\newcolumntype{A}{>{\PBS\raggedright\small\hspace{0pt}}X}
				\newcolumntype{L}[1]{>{\PBS\raggedright\small\hspace{0pt}}p{#1}}
				\newcolumntype{R}[1]{>{\PBS\raggedleft\small\hspace{0pt}}p{#1}}
				\newcolumntype{C}[1]{>{\PBS\centering\small\hspace{0pt}}p{#1}}
				
				\makeindex
				
				\title{TD Tri}	
			\date{}
			\author{\scriptsize{}}
			\definecolor{light-gray}{gray}{0.8}
			\renewcommand{\headrulewidth}{0pt}
			\fancyhead[L]{
				\footnotesize\textsc{Haute \'Ecole de Bruxelles}\\
	    			\footnotesize\textsc{\'Ecole Sup\'erieure d'Informatique}
			}
			\fancyhead[R]{
				\footnotesize{Bachelor en Informatique}\\
				\footnotesize{Laboratoires Java} - 
			\footnotesize{1\`ere ann\'ee}}
				\fancyfoot[L]{ }
				\fancyfoot[C]{}
				\fancyfoot[R]{\scriptsize{\textcolor{gray}{version 2014-2015 (\today)}}}
				\pagestyle{plain}
				\reversemarginpar
				\usepackage{rotating}						
				\begin{document}
					\begin{textblock}{9}(2,3.2)
						\includegraphics[width=2cm]{../../../_templates/java/icons/logo-esi}
					\end{textblock}
				
				
				
				
				%\maketitle
				\cadretitre{TD1}{TD Tri}
				\thispagestyle{fancy}
        \marginpar{\begin{sideways}
            \begin{minipage}[t]{1cm}
            \begin{tiny}
            \includegraphics[width=1\linewidth,height=1\textheight,keepaspectratio=true]{../../../_templates/java/icons/cc-gris.jpg}
			\end{tiny}
			\end{minipage}
            \begin{minipage}[b]{19cm}
            \begin{tiny}
            \textcolor{gray}{Distribué sous licence Creative Commons Paternité - Partage à l'Identique 2.0 Belgique 
            (\texttt{http://creativecommons.org/licenses/by-sa/2.0/be/})
			\vspace{-1em}
			\\Les autorisations au-delà du champ de cette licence peuvent être obtenues à 
			\texttt{http://www.heb.be/esi}
			- \texttt{mcodutti@heb.be}
			}\end{tiny}
			\end{minipage}
        \end{sideways}}
            \begin{abstract}
			Dans ce chapitre nous voyons quelques algorithmes simples pour trier un ensemble d'informations : 
			recherche des maxima, tri par insertion et tri bulle dans un tableau. Des algorithmes
			plus efficaces seront vus en deuxi\`eme ann\'ee.
    
            \par
        \end{abstract}
				\vspace{-2em}\tableofcontents
				\pagestyle{plain}
            \clearpage
            \fancyhead[L,C,R]{}
            \fancyfoot[L,C]{}
            \fancyfoot[R]{ \scriptsize{\textcolor{gray}{
				TDTri - page \thepage}}}
				\thispagestyle{fancy}
				\pagestyle{fancy}
	   
            \section{Motivation}
        La recherche efficace d'information implique un tri pr\'ealable de celle-ci. En effet, si les
        donn\'ees ne sont pas class\'ees ou tri\'ees, le seul algorithme possible reviendrait \`a parcourir
        enti\`erement l'ensemble des informations. Pour exemple, il suffit d'imaginer un dictionnaire
        dans lequel les mots seraient m\'elang\'es de fa\c con al\'eatoire au lieu d'\^etre class\'es par ordre
        alphab\'etique. Pour trouver le moindre mot dans ce dictionnaire, il faudrait \`a chaque fois le
        parcourir enti\`erement ! Il est clair que le classement pr\'ealable (ordre alphab\'etique) acc\'el\`ere
        grandement la recherche.
		
            \par
        
		  Ainsi, recherche et tri sont \'etroitement li\'es, et la fa\c con dont les informations sont tri\'ees
      conditionne bien entendu la fa\c con de rechercher l'information (cf. algorithme de recherche
      dichotomique). Pour exemple, prenons cette fois-ci un dictionnaire des mots crois\'es dans
      lequel les mots sont d'abord regroup\'es selon leur longueur et ensuite par ordre alphab\'etique.
      La fa\c con de rechercher un mot dans ce dictionnaire est bien s\^ur diff\'erente de la recherche
      dans un dictionnaire usuel.
		
            \par
        
		  Le probl\`eme central est donc le tri des informations. Celui-ci a pour but d'organiser un
      ensemble d'informations qui ne l'est pas \`a priori. On peut distinguer trois grands cas de
      figure :
		
            \par
        
			
		\subparagraph{donn\'ees \guillemotleft  brutes \guillemotright } 
		
					\textcolor{white}{.} \par
				
        D'abord les situations impliquant le classement total d'un ensemble de donn\'ees \guillemotleft  brutes \guillemotright , 
        c'est-\`a-dire compl\`etement d\'esordonn\'ees. Prenons pour exemple les feuilles r\'ecolt\'ees 
        en vrac \`a l'issue d'un examen ; il y a peu de chances que celles-ci soient remises \`a 
        l'examinateur de mani\`ere ordonn\'ee ; celui-ci devra donc proc\'eder au tri de l'ensemble 
        des copies, par exemple par ordre alphab\'etique des noms des \'etudiants, ou par num\'ero
        de groupe etc.
      
            \par
        
			
		\subparagraph{ajout dans un ensemble tri\'e} 
		
					\textcolor{white}{.} \par
				
		    Ensuite les situations o\`u on s'arrange pour ne jamais devoir trier la totalit\'e des \'el\'ements
        d'un ensemble, qui resterait cependant \`a tout moment ordonn\'e. Imaginons le cas d'une
        biblioth\`eque dont les livres sont rang\'es par ordre alphab\'etique des auteurs : \`a l'achat
        d'un nouveau livre, ou au retour de pr\^et d'un livre, celui-ci est imm\'ediatement rang\'e
        \`a la bonne place. Ainsi, l'ordre global de la biblioth\`eque est maintenu par la r\'ep\'etition
        d'une seule op\'eration \'el\'ementaire consistant \`a ins\'erer \`a la bonne place un livre parmi
        la collection. C'est la situation que nous consid\'ererions dans le cas d'une structure o\`u
        les \'el\'ements sont ordonn\'es.
      
            \par
        
			
		\subparagraph{re-trier} 
		
					\textcolor{white}{.} \par
				
        Enfin, les situations qui consistent \`a devoir re-trier des donn\'ees pr\'ealablement ordonn\'ees 
        sur un autre crit\`ere. Prenons l'exemple d'un paquet de copies d'examen d\'ej\`a tri\'ees
        sur l'ordre alphab\'etique des noms des \'etudiants, et qu'on veut re-trier cette fois-ci sur
        les num\'eros de groupe. Il est clair qu'une m\'ethode efficace veillera \`a conserver l'ordre
        alphab\'etique d\'ej\`a pr\'esent dans la premi\`ere situation afin que les copies apparaissent
        dans cet ordre dans chacun des groupes.
      
            \par
        
        Le dernier cas illustre un classement sur une cl\'e complexe (ou compos\'ee) impliquant la
        comparaison de plusieurs champs d'une m\^eme structure : le premier classement se fait sur le
        num\'ero de groupe, et \`a num\'ero de groupe \'egal, l'ordre se d\'epartage sur le nom de l'\'etudiant.
        On dira de cet ensemble qu'il est class\'e en \textbf{majeur}
        sur le num\'ero de groupe et en \textbf{mineur} sur le nom d'\'etudiant.
      
            \par
        
        Notons que certains tris sont dits \textbf{stables} 
        parce qu'en cas de tri sur une nouvelle cl\'e, l'ordre
        de la cl\'e pr\'ec\'edente est pr\'eserv\'e pour des valeurs identiques de la nouvelle cl\'e, ce qui \'evite
        de faire des comparaisons sur les deux champs \`a la fois. Les m\'ethodes nomm\'ees tri par
        insertion, tri bulle et tri par recherche de minima successifs (que nous allons aborder)
        sont stables.
      
            \par
        \subsection{Tri par insertion}
		    Cette m\'ethode de tri repose sur le principe d'insertion de valeurs dans un tableau ordonn\'e.
      
            \par
        
        Le tableau \`a trier sera \`a chaque \'etape subdivis\'e en deux sous-tableaux : le premier cadr\'e \`a
        gauche contiendra des \'el\'ements d\'ej\`a ordonn\'es, et le second, cadr\'e \`a droite, ceux qu'il reste
        \`a ins\'erer dans le sous-tableau tri\'e. Celui-ci verra sa taille s'accroitre au fur et \`a mesure des
        insertions, tandis que celle du sous-tableau des \'el\'ements non tri\'es diminuera progressivement.
        Au d\'epart de l'algorithme, le sous-tableau tri\'e est le premier \'el\'ement du tableau. Comme il
        ne poss\`ede qu'un seul \'el\'ement, ce sous-tableau est donc bien ordonn\'e ! Chaque \'etape consiste
        ensuite \`a prendre le premier \'el\'ement du sous-tableau non tri\'e et \`a l'ins\'erer \`a la bonne place
        dans le sous-tableau tri\'e.
		  
            \par
        \begin{verbatim}
/// Trie le tableau reçu en paramètre (via un tri par insertion).
module triInsertion(tab↓↑ : tableau de n entiers)
    i, j, valAInsérer : entiers
    pour i de 1 à n-1 faire
      valAInsérer ← tab[i]
      // recherche de l’endroit où insérer valAInsérer dans le
      // sous-tableau trié et décalage simultané des éléments
      j ← i-1
      tant que j ≥ 0 ET valAInsérer < tab[j] faire
        tab[j+1] ← tab[j]
        j ← j–1
      fin tant que
      tab[j+1] ← valAInsérer
    fin pour
fin module
\end{verbatim}\subsection{Tri par s\'election des minima successifs}
		    Dans ce tri, on recherche \`a chaque \'etape la plus petite valeur de l'ensemble non encore tri\'e
        et on peut la placer imm\'ediatement \`a sa position d\'efinitive.
      
            \par
        \begin{verbatim}
// Trie le tableau reçu en paramètre (via un tri par sélection des minima successifs).
module triSélectionMinimaSuccessifs(tab↓↑ : tableau de n entiers)
    i, indiceMin : entier
    pour i de 0 à n – 2 faire // i correspond à l’étape de l’algorithme
      indiceMin ← positionMin( tab, i, n-1 )
      swap( tab[i], tab[indiceMin] )
    fin pour
fin module


// Échange le contenu de 2 variables.
module swap(a↓↑, b↓↑ : entiers)
    aux : entiers
    aux ← a
    a ← b
    b ← aux
fin module


// Retourne l’indice du minimum entre les indices début et fin du tableau reçu.
module positionMin(tab↓ : tableau de n entiers, début, fin : entiers) → entier
    indiceMin, i : entiers
    indiceMin ← début
    pour i de début+1 à fin faire
      si tab[i] < tab[indiceMin] alors
        indiceMin ← i
      fin si
    fin pour
    retourner indiceMin
fin module

      \end{verbatim}\subsection{Tri bulle}
		    Il s'agit d'un tri par permutations ayant pour but d'amener \`a chaque \'etape \`a la \guillemotleft  surface \guillemotright 
        du sous-tableau non tri\'e (on entend par l\`a l'\'el\'ement d'indice minimum) la valeur la plus
        petite, appel\'ee la bulle. La caract\'eristique de cette m\'ethode est que les comparaisons ne se
        font qu'entre \'el\'ements cons\'ecutifs du tableau.
		  
            \par
        \begin{verbatim}
// Trie le tableau reçu en paramètre (via un tri bulle).
module triBulle(tab↓↑ : tableau de n entiers)
    indiceBulle, i : entiers
    pour indiceBulle de 1 à n faire
      pour i de n – 2 à indiceBulle par – 1 faire
        si tab[i] > tab[i + 1] alors
          swap( tab[i], tab[i + 1] )
        fin si
      fin pour
    fin pour
fin module


// Échange le contenu de 2 variables.
module swap(a↓↑, b↓↑ : entiers)
    aux : entiers
    aux ← a
    a ← b
    b ← aux
fin module
\end{verbatim}\subsection{Recherche dichotomique}
		    La recherche dichotomique a pour essence de r\'eduire \`a chaque \'etape la taille de l'ensemble
        de recherche de moiti\'e, jusqu'\`a ce qu'il ne reste qu'un seul \'el\'ement dont la valeur devrait
        \^etre celle recherch\'ee, sauf bien entendu en cas d'inexistence de cette valeur dans l'ensemble
        de d\'epart.
		  
            \par
        
		    Soit val la valeur recherch\'ee dans une zone d\'elimit\'ee par les indices indiceGauche et indiceDroit. 
		    On commence par d\'eterminer l'\'el\'ement m\'edian, c'est-\`a-dire celui qui se trouve \guillemotleft  au
        milieu \guillemotright  de la zone de recherche ; son indice sera d\'etermin\'e par la formule
        \,\verb|indiceMédian ← (indiceGauche + indiceDroit) DIV 2|\,
            \par
        
        N.B. : cet \'el\'ement m\'edian n'est pas tout \`a fait au milieu dans le cas d'une zone contenant un
        nombre pair d'\'el\'ements. On compare alors val avec la valeur de cet \'el\'ement m\'edian ; il est
        possible qu'on ait trouv\'e la valeur cherch\'ee ; sinon, on partage la zone de recherche en deux
        parties : une qui ne contient certainement pas la valeur cherch\'ee et une qui pourrait la
        contenir. C'est cette deuxi\`eme partie qui devient la nouvelle zone de recherche. On r\'eit\`ere
        le processus jusqu'\`a ce que la valeur cherch\'ee soit trouv\'ee ou que la zone de recherche soit
        r\'eduite \`a sa plus simple expression, c'est-\`a-dire un seul \'el\'ement.
		  
            \par
        \begin{verbatim}

module rechercheDichotomique(tab↓↑ : tableau de n T, valeur↓ : T, pos↑ : entier) → booléen
    indiceDroit, indiceGauche, indiceMédian : entiers
    candidat : T
    trouvé : booléen
    indiceGauche ← 0
    indiceDroit ← n-1
    trouvé ← faux
    tant que NON trouvé ET indiceGauche ≤ indiceDroit faire
      indiceMédian ← (indiceGauche + indiceDroit) DIV 2
      candidat ← tab[indiceMédian]
      si candidat = valeur alors 
        trouvé ← vrai
      sinon si candidat < valeur alors
        indiceGauche ← indiceMédian + 1 // on garde la partie droite
      sinon 
        candidat > valeur: indiceDroit ← indiceMédian – 1 // on garde la partie gauche
      fin si
    fin tant que
    
    si trouvé alors
      pos ← indiceMédian
    sinon
      pos ← indiceGauche
      // dans le cas où la valeur n’est pas trouvée,
      // on vérifiera que indiceGauche donne la valeur où elle pourrait être insérée.
    fin si
    retourner trouvé
fin module
\end{verbatim}\section{Exercices}
				Maintenant, mettons tout \c ca en pratique.
      
            \par
        \subsection{En Java...}
        N'oubliez pas nos quelques conseils pour vous guider dans la r\'esolution de tels probl\`emes :
        
					\begin{itemize}
				
			\item il convient d'abord de bien comprendre le probl\`eme pos\'e ; assurez-vous qu'il est parfaitement sp\'ecifi\'e ;
			\item r\'esolvez le probl\`eme via quelques exemples pr\'ecis ;
			\item mettez en \'evidence les variables \textbf{\guillemotleft  donn\'ees \guillemotright }, les variables \textbf{\guillemotleft  r\'esultats \guillemotright } et les variables de travail ;
			\item n'h\'esitez pas \`a faire une \'ebauche de r\'esolution en fran\c cais avant d'\'elaborer l'algorithme d\'efinitif pseudo-cod\'e ;
			\item d\'eclarez ensuite les variables (et leur type) qui interviennent dans chaque module ; les noms des variables risquant de ne pas \^etre suffisamment explicites.
			\item \'Ecrivez la partie algorithmique \textbf{AVANT} de vous lancer dans la programmation en Java.
			\item Demandez-vous si vous avez besoin de parcourir tout le tableau ou de sortir pr\'ematur\'ement (si on a trouv\'e ce qu'on cherche par exemple).
			\item Pour la partie Java, dessinez l'arborescence des fichiers. 
			\item \textbf{\'Ecrivez le plan de tests en \'ecrivant l'algorithme. Codez les tests apr\`es avoir \'ecrit le code Java.}
					\end{itemize}
				
            \par
        
        \'Ecrivez en Java les algorithmes de tri vus ici en affichant le tableau \`a chaque \'etape.
      
            \par
        
			
		\subparagraph{S\'election des maxima} 
		
					\textcolor{white}{.} \par
				
        D\'eveloppez un algorithme similaire consistant \`a trier un tableau par ordre croissant par
        s\'election des maxima successifs. Le sous-tableau tri\'e apparaitra donc \`a droite du tableau,
        et les maxima s\'electionn\'es seront \`a chaque \'etape positionn\'es \`a droite du sous-tableau non tri\'e.
      
            \par
        
			
		\subparagraph{Maxima et minima} 
		
					\textcolor{white}{.} \par
				
        D\'eveloppez un algorithme combinant les deux recherches. \`A chaque \'etape, on s\'electionne
        donc le minimum et le maximum du sous-tableau restant \`a trier et on les positionnera \`a
        l'endroit ad hoc. Cette m\'ethode apporte-t-elle une am\'elioration en temps ou en simplicit\'e
        aux deux algorithmes de base ?
      
            \par
        
			
		\subparagraph{Am\'elioration du tri bulle} 
		
					\textcolor{white}{.} \par
				
        \'Ecrivez une am\'elioration du tri bulle consistant \`a m\'emoriser \`a chaque \'etape l'indice de la
        derni\`ere permutation, celui-ci d\'elimitant en fait la v\'eritable taille du sous-tableau tri\'e \`a
        l'issue d'une \'etape. En lieu et place de la boucle pour n'incr\'ementant l'indice bulle que de 1
        \`a la fois, vous \'ecrirez une boucle \verb@tant que@ 
        \`a l'issue de laquelle \verb@indiceBulle@ prendra la valeur de
        l'\textit{indice de derni\`ere permutation + 1}.
      
            \par
        
			
		\subparagraph{Tri shaker} 
		
					\textcolor{white}{.} \par
				
        L'am\'elioration pr\'ec\'edente est issue de l'observation du sous-tableau d\'ej\`a tri\'e en d\'ebut du
        tableau initial. On peut de m\^eme \'etudier la possibilit\'e d'avoir un sous-tableau, tri\'e \'egale-
        ment, mais cadr\'e \`a droite dans le tableau \`a trier. Cette sym\'etrie sugg\`ere une am\'elioration
        suppl\'ementaire qui consiste \`a changer de sens \`a la fin d'un parcours pour entamer le parcours
        suivant. Lorsqu'on change de sens, on am\`enera l'\'el\'ement le plus grand (qu'on peut nommer
        le \guillemotleft  plomb \guillemotright ) au fond du tableau non tri\'e. L'association des deux m\'ethodes donne ce qu'on
        appelle le tri shaker, dont le but est de restreindre le sous-tableau non tri\'e en augmentant
        sa borne inf\'erieure et en diminuant sa borne sup\'erieure. \'Ecrivez l'algorithme qui r\'ealise cette
        m\'ethode de tri.
      
            \par
        
				\end{document}
			